% ==============================================================================
% SLIDE 6: SEISMIC RESILIENCE & MODERN CODES
% ==============================================================================
\begin{frame}[label=seismic]{Seismic Resilience: Traditional Craft to Modern Code}
  \begin{columns}[T, onlytextwidth]
    % Left Column: Technical Diagram
    \begin{column}{0.48\textwidth}
      \slidephoto{akrotiri-wooden-bracing.jpg}{Bronze Age anti-seismic engineering: Timber-laced masonry (\textit{imantosis}) dissipating lateral shear at Akrotiri ($\sim$1600 BC)}
    \end{column}

    % Right Column: Bullet Points
    \begin{column}{0.48\textwidth}
      \textbf{\color{aegeanblue}Evolution of Seismic Protection:}
      \vspace{0.1em}
      \begin{itemize}
          \setlength{\itemsep}{0.08em}
        \item \textbf{Antiquity Precedents (Akrotiri $\sim$1600 BC):} Bronze Age builders embedded timber lacing into stone walls to absorb tension.
        \item \textbf{The 1956 Amorgos Disaster (M 7.7):} Widespread destruction revealed the fatal flaw of unreinforced stone tearing at corners.
        \item \textbf{Concrete Ring Beams (\textit{Chainage}):} Modern codes mandate continuous reinforced concrete bands at foundation, floor, and roof.
        \item \textbf{Monolithic Rigid-Box Behavior:} Binds disparate walls into a unified 3D box, preventing out-of-plane detachment.
      \end{itemize}
    \end{column}
  \end{columns}
\end{frame}
